% Chapter 1: expert onboarding, mathematical scope, and the evidentiary contract.
% Research baseline etblink/math@d4870eb1ac58, FCP historical receipts recorded separately.
\chapter{Scope, Physical Context, and the Research Problem}
\label{ch:scope}

\section{The object of study is a quantum system, not an assumed interpretation}

The subject of this volume is a specified supersymmetric matrix quantum
mechanical system: the \emph{relative} $N=2$ case of the Banks--Fischler--Shenker--Susskind
(BFSS) family. Its fundamental variables are nine traceless Hermitian
$2\times2$ matrix coordinates and sixteen matrix-valued fermionic
components. The relative adjective indicates that the decoupled trace or
center-of-mass sector is not part of the system analyzed here. The
original source fixes a concrete choice of color basis, Clifford matrices,
fermionic CAR realization, supercharge normalization, and physical gauge
action. We will not silently replace any of these data with a more convenient
isomorphic model when a theorem purports to concern the original system.

This fixed-$N$ quantum mechanics is mathematically meaningful whether or not
one assumes a particular high-energy interpretation. The original BFSS
proposal relates a large-$N$ matrix-model limit to eleven-dimensional
M-theory. That conjectural relationship is the historical motivation for the
name; it is \emph{not} a hypothesis of our source definitions, of the
physical Hilbert-space construction, or of the present spectral program.
Nine spatial matrix labels do not amount to eleven explicitly modeled
spacetime coordinates. Conversely, the 27 real coefficients in the
relative $N=2$ bosonic configuration space are matrix components, not a
27-dimensional physical spacetime. This distinction will be maintained
throughout the book.

\begin{definition}[The fixed mathematical research target]
\label{def:research-target}
The \emph{original relative $\mathrm{SU}(2)$ system} means the source
$N=2$ specialization of \texttt{OAI.BFSSQuantum.AlgebraData} together with
its genuine gauge representation, original smooth core, sixteen original
supercharges, and the associated physical closed charge/Hamiltonian
construction, with the source's exact representation and normalization data.
A statement about an auxiliary Fock algebra or a different matrix model is
not by this definition a statement about the original physical Hamiltonian.
\end{definition}

The underlying carrier will be introduced precisely in
Chapter~\ref{ch:carriers}. In advance, its schematic form is
\begin{equation}\label{eq:intro-full-hilbert}
 \mathscr H=L^2(\mathbb R^{27},dx;\mathcal F),\qquad
 \dim_{\mathbb C}\mathcal F=2^{24}.
\end{equation}
The physical Hilbert space $\mathscr H_{\mathrm{phys}}$ is the closed
subspace satisfying the \emph{actual simultaneous} $\mathrm{SU}(2)$
color-gauge invariance. It is not the tensor product of a purported
three-family flavor space and an unrelated spinor Hilbert space. The
three real color directions form the adjoint gauge representation.

\section{Why the operator and its domain must precede the spectral question}

The compact formula for a supersymmetric Hamiltonian conceals its most
important analytic issue: a differential expression is not yet a
self-adjoint unbounded operator. At the smooth-core level the original
sixteen charge maps $Q_\alpha$ determine the nonnegative real form
\begin{equation}\label{eq:intro-core-form}
 q_{\mathrm{core}}(\psi)
 =\frac1{16}\sum_{\alpha=0}^{15}
   \bigl\|\iota(Q_\alpha\psi)\bigr\|_{\mathscr H}^{2},
 \qquad \psi\in\mathscr D_{\mathrm{phys}},
\end{equation}
where $\iota$ denotes the original test-function-to-$L^2$ map and
$\mathscr D_{\mathrm{phys}}$ is the source physical smooth core. The
factor $1/16$ is not an adjustable unit convention in a source-exact
claim. It is the one occurring in the formalized source
\texttt{AlgebraData.coreForm}.

The historical physical-domain program proved density, closedness and
natural-domain statements leading to the native nonnegative operator
\begin{equation}\label{eq:intro-physical-hamiltonian}
 H_{\mathrm{phys}}=\frac1{16}\,T^\dagger T,
 \qquad
 \operatorname{Dom}(H_{\mathrm{phys}})
 =\{\psi\in\operatorname{Dom}(T):
       T\psi\in\operatorname{Dom}(T^\dagger)\}.
\end{equation}
Here $T$ denotes the \emph{closed original physical sixteen-charge
column} in the historically qualified construction; the formula does
not assert that an arbitrary formal core column has the same domain.
The BA--BH sequence and its exact historical receipts are explained
in Chapter~\ref{ch:hamiltonian} and the historical appendix. Importing
those Lean files into a new tree without a separately qualified cold
consumer does not create a new integration-level acceptance certificate.

\section{The spectral questions: point spectrum, compact resolvent, and a gap}

Three superficially similar questions must be separated. A
\emph{normalizable eigenstate} is a nonzero vector
$\psi\in\operatorname{Dom}(H_{\mathrm{phys}})$ satisfying
\begin{equation}\label{eq:intro-eigenvalue}
 H_{\mathrm{phys}}\psi=E\psi,\qquad
 0<\|\psi\|_{\mathscr H}<\infty.
\end{equation}
An eigenvalue is \emph{positive} when $E>0$. Its existence is a
point-spectrum assertion and does not require the entire physical
spectrum to be discrete. A \emph{compact-resolvent theorem on a
reducing sector} is stronger in one direction and weaker in another:
when the sector is infinite-dimensional, the restricted operator has
an unbounded discrete sequence of eigenvalues, but the restriction
need not describe the rest of the full physical spectrum. A
\emph{positive spectral gap of the entire physical Hamiltonian} is
a different claim about the global low-energy spectrum; it cannot be
inferred from a positive sector eigenpair.

Our primary prospective theorem concerns positive point spectrum in
a rigorously identified original physical sector. We do \emph{not}
assert a mass gap or positive-energy eigenstate as an established
conclusion. Classical flat directions of the quartic bosonic potential
make careless global confinement claims particularly suspect. The
source exact energy identity is therefore a starting point, not a
substitute for an estimate excluding escape to infinity.

\begin{remark}[The meaning of ``original'']
We distinguish the original physical $\mathcal F$-valued Hilbert
space~\eqref{eq:intro-full-hilbert} from the auxiliary
$\operatorname{End}_{\mathbb C}(\mathcal F)$-valued Hilbert space used
in a conditional oscillator argument. Their representations and
highest-weight bounds differ. A theorem proved for the second space
cannot be transferred to the first by changing notation.
\end{remark}

\section{Three mathematical proof chains and the unproved joins}

The research is easier to understand as three interacting chains
than as a chronological list of CI jobs.

\begin{description}[leftmargin=1.4em,style=nextline]
 \item[The physical-operator chain.]
 Original supercharges and gauge action lead to the physical smooth
 core, density of its charge domain, a closed charge column, the
 normalized adjoint-square Hamiltonian, and historical source-native
 operator identification. Chapters~\ref{ch:charges}--\ref{ch:hamiltonian}
 distinguish these statements and their exact domain requirements.
 \item[The spatial-representation chain.]
 Source $\operatorname{Spin}(9)$, its original nine-vector and
 48-Majorana actions, the exact ninth-axis covering, Clifford--Fock
 operator automorphisms and genuine gauge-label compatibility have
 substantial formal proof support. A continuous source-exact
 \emph{Fock-vector} action preserving the full physical charge column
 and its closed domain is a further mathematical obligation.
 Chapters~\ref{ch:spin-geometry}--\ref{ch:covariance}
 make this boundary explicit.
 \item[The energy-and-compactness chain.]
 The original massless core energy decomposition and the abstract
 first-radial-moment compactness theorems are qualified source results.
 The necessary high-type \emph{physical-sector coercivity}, source
 fermionic field estimate, and closed-form-domain transfer remain
 open. Chapters~\ref{ch:energy}--\ref{ch:spectral} reconstruct the
 conditional bridge without treating it as an accomplished theorem.
\end{description}

The central missing join can be written schematically as
\begin{equation}\label{eq:intro-join}
\begin{gathered}
 \text{original physical Hamiltonian and source-compatible sector}\\
 \Big\Downarrow\quad\text{\emph{open}: coercivity and closed-domain transfer}\\
 \text{first-radial-moment control and verified abstract compactness}\\
 \Big\Downarrow\quad\text{\emph{conditional}: closed-form spectral theorem}\\
 \text{positive sector eigenvalues}.
\end{gathered}
\end{equation}
The arrow marked \emph{open} is the current mathematical research
frontier. The right-hand implication is a standard theorem of closed
quadratic forms, \emph{conditional on the precise hypotheses}.

\section{Verified, conditional, and exploratory statements are different data}

This monograph is simultaneously a mathematical exposition and a
reproducible record of what the formal project has actually accepted.
We use the following \emph{evidence statuses}, indexed in
\texttt{THEOREM\_ATLAS.json}:
\begin{description}[leftmargin=1.4em,style=nextline]
 \item[Original source definition.] A declaration specifying an exact
 carrier, operator or normalization; a definition alone proves no
 spectral consequence.
 \item[Historically strict-qualified.] A named theorem and its
 independent consumer/axiom evidence at an exact tested historical
 commit, often in the FCP record. It is not automatically a green
 result for an altered copy on our current branch.
 \item[Integrated strict-qualified.] A source theorem independently
 cold-compiled and checked in the \texttt{etblink/math} research
 lineage at a specified tested commit.
 \item[Source present or qualification pending.] A declaration may
 exist in the tree while the required exact acceptance is unverified,
 queued, or failed. The distinction is substantive.
 \item[Conditional analytical argument.] A conventional derivation
 from explicit premises that are not yet bound to the exact original
 physical system. It may be mathematically illuminating but has not
 closed the source-specific physical theorem.
 \item[Open proof obligation.] A statement selected as a research
 target, rather than an established result.
\end{description}
A green infrastructure check is evidence about a particular source
and consumer contract, not a substitute for an independent
mathematical explanation. Conversely, a failure due to a generated
manifest, file path, or deprecated Lean tactic is not in itself a
counterexample to a mathematical statement. Both successes and
failures are preserved at their historical commit and scope.

\section{What the model cannot yet tell us about the universe}

The source Spin(9) representation, original color triplet, Clifford
algebra and gauge symmetry are mathematically real. They do not by
themselves establish four-dimensional chiral matter, three particle
generations, the Standard Model gauge and charge assignments, spacetime
emergence, or a theory of quantum gravity. In particular, three adjoint
color components are not three gauge-invariant flavor labels. The
putative relation between a real Spin(9) sixteen and a complex chiral
Spin(10) sixteen is representation-theoretic motivation, not an
identification of their physical roles.

The later exploratory chapters examine candidate-neutral connections
to Hopf geometry, gauge descent, protected topology, category theory,
and particle unification. Every proposal must specify a mathematical
map on the \emph{actual} source carriers, preserve gauge equivalence
and physical operator domains, and lead eventually to new testable
physical content if it is to be promoted beyond structural analogy.
Unsuccessful constructions are useful negative evidence and are
recorded without disguising their failures.

\section{How to read and extend the book}

Chapters~\ref{ch:carriers}--\ref{ch:hamiltonian} first reconstruct the
system's source definition and genuine physical Hamiltonian.
Chapters~\ref{ch:spin-geometry}--\ref{ch:covariance} develop
spatial symmetry and its strict representation boundaries.
Chapters~\ref{ch:energy}--\ref{ch:spectral} explain the energy,
normalized auxiliary observables, conditional oscillator exclusion,
and the compactness-to-spectrum argument. The research-frontier and
prospective-unification chapters then separate the remaining physical
proof obligations from genuinely new hypotheses.

The main exposition follows \emph{logical dependency order}. The
historical appendix follows \emph{chronological order}, from the
October FCP charge/gauge/domain results to their later integration and
the November Spin(9), Fock and spectral interfaces. The theorem atlas
provides the exact source paths, immutable Git blobs and strict-run
receipts; its dependency relationships are editorial and must not be
mistaken for Lean's import graph. The appendices distinguish source
notations, normalization choices and negative knowledge. A new
researcher should be able to locate both a mathematical proof and
its formal acceptance record without confusing them.

\formalrecord{Research baseline: \texttt{etblink/math@d4870eb1ac58};
original carriers and charges in \texttt{Core.lean}; historical FCP
physical-domain BA--BH sources and receipts in
\texttt{HISTORICAL\_LINEAGE.md}; integrated spatial, Fock and spectral
qualification in companion \texttt{THEOREM\_ATLAS.json}. This chapter
introduces no new Lean theorem. The historical BFSS large-$N$ conjecture
is Banks--Fischler--Shenker--Susskind (1997),
\href{https://doi.org/10.1103/PhysRevD.55.5112}{doi:10.1103/PhysRevD.55.5112}.}
